\documentclass[12pt]{article}
\usepackage{graphicx}
\usepackage[margin=1in]{geometry}
\usepackage{caption}
\usepackage{xcolor}
\usepackage{float}
\usepackage{amsmath}
\usepackage{array}
\usepackage{multicol}
\usepackage{hyperref}
\usepackage{url}
\usepackage{listings}
\usepackage{courier}

\lstset{
  basicstyle=\ttfamily\footnotesize,
  backgroundcolor=\color{gray!10},
  frame=single,
  breaklines=true,
  captionpos=b
}

\setlength{\columnsep}{1cm}

\begin{document}

% Title Block
\begin{figure}[H]
    \begin{minipage}{0.45\textwidth}
        \includegraphics[width=\textwidth]{iiitb_comet.jpeg} % Replace with your image
    \end{minipage} \hfill
    \begin{minipage}{0.45\textwidth}
        \textbf{Name: S. Rithrisha Reddy} \\
        \textbf{Batch: COMETFWC028} \\
        \textbf{Date: 27 May 2025}
    \end{minipage}
\end{figure}

\begin{center}
    {\LARGE \textbf{\textcolor{cyan}{GATE 2009, ECE Question Number 36}}}
\end{center}

\begin{multicols}{2}

\vspace{1em}
\noindent\textbf{Abstract} \\
This project implements and verifies a logic equation derived from a GATE 2009 ECE question using AVR-GCC and physical hardware (ATmega328P). Logical analysis determines the value of variable \(Z\), and the expression is validated using an LED output.

\vspace{1em}
\noindent\textbf{1. Question and Simplification}
\begin{itemize}
    \item Given: \( X = 1 \)
    \item Expression: \( (X + Z)(Y + (Z + Y))(X + Z(X + Y)) = 1 \)
    \item Simplifies to: \( (1)(Y + Z)(1 + Z(1 + Y)) = 1 \)
    \item Solving: \( Z(1 + Y) = 0 \Rightarrow Z = 0 \)
\end{itemize}
\textbf{Correct Answer: (D) Z = 0}

\vspace{1em}
\noindent\textbf{2. Components}
\begin{table}[H]
\small
\centering
\begin{tabular}{|p{4.2cm}|c|}
\hline
\textbf{Component} & \textbf{Qty} \\
\hline
ATmega328P / AVR Board & 1 \\
Push Buttons (X, Y, Z) & 3 \\
LED for Output & 1 \\
Resistors (220$\Omega$, 10k$\Omega$) & 4 \\
Breadboard & 1 \\
Jumper Wires & 10 \\
Laptop with AVR-GCC & 1 \\
\hline
\end{tabular}
\caption*{Table 1: Components used}
\end{table}

\vspace{1em}
\noindent\textbf{3. Setup and Connections}
\begin{itemize}
    \item Connect buttons to PD0 (X), PD1 (Y), PD2 (Z).
    \item Connect LED to PB0 with 220$\Omega$ resistor.
    \item Use 10k$\Omega$ pull-down resistors for each button.
    \item Ensure common GND and VCC rails.
\end{itemize}

\vspace{1em}
\noindent\textbf{4. Logic Summary}
\begin{itemize}
    \item Inputs: X, Y, Z (via push buttons)
    \item Evaluated Expression: \( (X + Z)(Y + Z)(1 + Z(1 + Y)) \)
    \item Output (LED) glows only if the expression evaluates to 1
\end{itemize}

\vspace{1em}
\noindent\textbf{5. Pin Mapping}
\begin{itemize}
    \item \textbf{X} – PD0 (Input)
    \item \textbf{Y} – PD1 (Input)
    \item \textbf{Z} – PD2 (Input)
    \item \textbf{LED} – PB0 (Output)
\end{itemize}

\end{multicols}

\section*{6. Verification Procedure}
\begin{enumerate}
    \item Set up the circuit on a breadboard using the components listed.
    \item Ensure pull-down resistors are connected properly for X, Y, and Z inputs.
    \item Burn the compiled AVR-GCC C code into the ATmega328P microcontroller.
    \item Toggle the buttons to test various input combinations.
    \item Observe LED state for each combination.
    \item Verify output matches expected logic from truth table.
\end{enumerate}

\section*{7. Analysis and Results}

\subsection*{7.1 Sample Truth Table}
\begin{table}[H]
\centering
\begin{tabular}{|c|c|c|c|}
\hline
X & Y & Z & Output \\
\hline
1 & 0 & 0 & 1 \\
1 & 1 & 0 & 1 \\
1 & 0 & 1 & 0 \\
1 & 1 & 1 & 0 \\
\hline
\end{tabular}
\caption*{Table 2: Logic validation for Z}
\end{table}

\subsection*{7.2 Circuit Implementation}
\begin{figure}[H]
\centering
\includegraphics[width=0.8\linewidth]{output1.jpeg}
\caption*{Figure: Practical Circuit Using AVR and Breadboard}
\end{figure}

\section*{8. Conclusion}
This project successfully demonstrated the evaluation of a logical expression from GATE 2009 using AVR hardware and C programming. The experimental LED output confirmed that the function outputs 1 only when \( Z = 0 \), validating the theoretical simplification.

\section*{9. References}
\begin{itemize}
    \item GitHub Repository: \url{https://github.com/Rithrishareddy/RITHRISHA_FWC/tree/main/Hardware}
    \item GATE 2009 ECE Question Paper
    \item AVR-GCC Documentation and Datasheets
\end{itemize}

\end{document}
